\subsection{SEMI-CONTINUOUS FUNCTIONS ON A BAIRE SPACE}

THEOREM 2. Let X be a Baire space and let \((f_{\alpha})\) be a family of lower semi-continuous real-valued functions on X such that, at every point x of X, the upper envelope \(\sup_{\alpha} f_{\alpha}(x)\) is finite. Then every non-empty open set in X contains a non-empty open set on which the family \((f_{\alpha})\) is uniformly bounded above.

This theorem may also be stated in the form that the set of points in the neighbourhood of which the family \((f_{\alpha})\) is uniformly bounded above is a dense open set.

Let \(f = \sup_{\alpha} f_{\alpha}\) be the upper envelope of the family \((f_{\alpha})\) . The function \(f\) is lower semi-continuous (Chapter IV, § 6, no. 2, Theorem 4) and finite at every point of \(\mathbf{X}\) . It is therefore enough to carry through the proof in the case where the family \((f_{\alpha})\) consists of a single function \(f\) . Let \(\mathbf{A}_n\) be the set of points \(x \in \mathbf{X}\) such that \(f(x) \leqslant n\) ; \(\mathbf{A}_n\) is closed (Chapter IV, § 6, no. 2, Proposition 1), and the assumptions imply that \(\mathbf{X}\) is the union of the sets \(\mathbf{A}_n\) ; hence at least one of the \(\mathbf{A}_n\) has an interior point, and therefore there is a non-empty open set on which \(f\) is bounded above (by an integer \(n\) ). If we apply this result to any non-empty open subset of \(\mathbf{X}\) (this subspace is also a Baire space by Proposition 3 of no. 3), we have the theorem.

In applications of this theorem it is generally the case that the \(f_{\alpha}\) are continuous on X.

Remark. The theorem may be false if we do not suppose that X is a Baire space. For example if, for each rational number p/q (p, q being coprime integers, q \textgreater{} 0) we put \(f(p/q) = q\) , we have a lower semicontinuous function f on the rational line Q which is finite at each point (cf.~Chapter IV, § 6, no. 2); but there is no non-empty open set in Q on which f is bounded above.

